% GENERATED FILE. Edit canonical lessons, then run npm run generate:books.
% canonical-source-hash: fnv1a64:c649692cb20963cc
% canonical-lessons: TA-C450-acal, TA-C450-poli-fake, TA-C450-atamparam, TA-C450-vilaivaci, TA-C450-nottu

\chapter{Real or Fake}
\label{ch:ta-real-or-fake}
\hlchaptermodality{450}

\begin{chapteropening}
\textbf{By the end of this chapter:} \emph{I can talk about originals and fakes, luxury, the cost of living and banknotes.}

Complete the last lesson of chapter 450: I can talk about originals and fakes, luxury, the cost of living and banknotes.
\end{chapteropening}

\section[\texorpdfstring{\ta{அசல்} (acal)}{அசல் (acal)}]{\ta{அசல்} (acal) — the original; genuine}

\label{lesson:TA-C450-acal}



\begin{quote}
Before the new one: say the Tamil for an investment, then the Tamil for to pawn.
\end{quote}

\subsection*{You'll want to know: \ta{அசல்}}
\textbf{\ta{அசல்}} — \emph{acal} — \textquotedblleft{}the original; genuine\textquotedblright{}.

\begin{quote}
\ta{இது} \ta{அசல்} \ta{சான்றிதழ்}, \ta{நகல்} \ta{அல்ல}.
\end{quote}

\begin{quote}
\emph{itu acal cāṉṟitaḻ, nakal alla.}
\end{quote}

\begin{quote}
This is the original certificate, not a copy.
\end{quote}

The opposite is \ta{நகல்}, a copy.

\begin{cousinweb}
\textbf{\ta{அசல்}} (\emph{acal}) is a loan from Arabic \emph{aṣl}, \textquotedblleft{}root, origin\textquotedblright{}, through Hindustani.
\end{cousinweb}

\begin{grammarlens}[title={What you've built}]
One more word for getting around.
\end{grammarlens}

\subsection*{Your turn}
\begin{itemize}
\raggedright
  \item \emph{Say it:} \emph{acal}
  \item \emph{Say it:} \emph{acal}, once more
  \item \emph{Say it:} \emph{itu acal cāṉṟitaḻ, nakal alla.} — This is the original certificate, not a copy.
  \item \emph{Recall:} say the Tamil for an investment, then the Tamil for to pawn, then say \emph{acal} again
\end{itemize}

\begin{tcolorbox}[breakable,colback=teal!4,colframe=teal!35!black,title={Before you move on}]
What does \ta{அசல்} mean? (\textquotedblleft{}The original; genuine\textquotedblright{}.) Say it once more.
\end{tcolorbox}
\section[\texorpdfstring{\ta{போலி} (pōli)}{போலி (pōli)}]{\ta{போலி} (pōli) — a fake; an imitation}

\label{lesson:TA-C450-poli-fake}



\begin{quote}
Before the new one: say the Tamil for to pawn, then the Tamil for the original.
\end{quote}

\subsection*{You'll want to know: \ta{போலி}}
\textbf{\ta{போலி}} — \emph{pōli} — \textquotedblleft{}a fake; an imitation\textquotedblright{}.

\begin{quote}
\ta{அந்தக்} \ta{கடையில்} \ta{போலிப்} \ta{பொருட்களை} \ta{விற்கிறார்கள்}.
\end{quote}

\begin{quote}
\emph{antak kaṭaiyil pōlip poruṭkaḷai viṟkiṟārkaḷ.}
\end{quote}

\begin{quote}
They sell fake goods in that shop.
\end{quote}

As a modifier: \ta{போலிப்} \ta{பொருள்}, a fake item.

\begin{cousinweb}
\textbf{\ta{போலி}} (\emph{pōli}) is native Dravidian (DEDR 4597), \textquotedblleft{}likeness, a pretence\textquotedblright{}, from the root of \ta{போல்}, \textquotedblleft{}like\textquotedblright{}.
\end{cousinweb}

\begin{grammarlens}[title={What you've built}]
One more word for getting around.
\end{grammarlens}

\subsection*{Your turn}
\begin{itemize}
\raggedright
  \item \emph{Say it:} \emph{pōli}
  \item \emph{Say it:} \emph{pōli}, once more
  \item \emph{Say it:} \emph{antak kaṭaiyil pōlip poruṭkaḷai viṟkiṟārkaḷ.} — They sell fake goods in that shop.
  \item \emph{Recall:} say the Tamil for to pawn, then the Tamil for the original, then say \emph{pōli} again
\end{itemize}

\begin{tcolorbox}[breakable,colback=teal!4,colframe=teal!35!black,title={Before you move on}]
What does \ta{போலி} mean? (\textquotedblleft{}A fake; an imitation\textquotedblright{}.) Say it once more.
\end{tcolorbox}
\section[\texorpdfstring{\ta{ஆடம்பரம்} (āṭamparam)}{ஆடம்பரம் (āṭamparam)}]{\ta{ஆடம்பரம்} (āṭamparam) — luxury; show, pomp}

\label{lesson:TA-C450-atamparam}



\begin{quote}
Before the new one: say the Tamil for the original, then the Tamil for a fake.
\end{quote}

\subsection*{You'll want to know: \ta{ஆடம்பரம்}}
\textbf{\ta{ஆடம்பரம்}} — \emph{āṭamparam} — \textquotedblleft{}luxury; show, pomp\textquotedblright{}.

\begin{quote}
\ta{அவர்கள்} \ta{திருமணத்தை} \ta{ஆடம்பரம்} \ta{இல்லாமல்} \ta{நடத்தினார்கள்}.
\end{quote}

\begin{quote}
\emph{avarkaḷ tirumaṇattai āṭamparam illāmal naṭattiṉārkaḷ.}
\end{quote}

\begin{quote}
They held the wedding without any show.
\end{quote}

\ta{ஆடம்பரமான}, luxurious, is the describing form.

\begin{cousinweb}
\textbf{\ta{ஆடம்பரம்}} (\emph{āṭamparam}) is borrowed from Sanskrit \emph{āḍambara}, \textquotedblleft{}a great noise, bombast\textquotedblright{} (Monier-Williams); in Tamil it came to mean show and pomp.
\end{cousinweb}

\begin{grammarlens}[title={What you've built}]
One more word for getting around.
\end{grammarlens}

\subsection*{Your turn}
\begin{itemize}
\raggedright
  \item \emph{Say it:} \emph{āṭamparam}
  \item \emph{Say it:} \emph{āṭamparam}, once more
  \item \emph{Say it:} \emph{avarkaḷ tirumaṇattai āṭamparam illāmal naṭattiṉārkaḷ.} — They held the wedding without any show.
  \item \emph{Recall:} say the Tamil for the original, then the Tamil for a fake, then say \emph{āṭamparam} again
\end{itemize}

\begin{tcolorbox}[breakable,colback=teal!4,colframe=teal!35!black,title={Before you move on}]
What does \ta{ஆடம்பரம்} mean? (\textquotedblleft{}Luxury; show, pomp\textquotedblright{}.) Say it once more.
\end{tcolorbox}
\section[\texorpdfstring{\ta{விலைவாசி} (vilaivāci)}{விலைவாசி (vilaivāci)}]{\ta{விலைவாசி} (vilaivāci) — the cost of living; prices}

\label{lesson:TA-C450-vilaivaci}



\begin{quote}
Before the new one: say the Tamil for a fake, then the Tamil for luxury.
\end{quote}

\subsection*{You'll want to know: \ta{விலைவாசி}}
\textbf{\ta{விலைவாசி}} — \emph{vilaivāci} — \textquotedblleft{}the cost of living; prices\textquotedblright{}.

\begin{quote}
\ta{இந்த} \ta{ஆண்டு} \ta{விலைவாசி} \ta{மிகவும்} \ta{ஏறிவிட்டது}.
\end{quote}

\begin{quote}
\emph{inta āṇṭu vilaivāci mikavum ēṟiviṭṭatu.}
\end{quote}

\begin{quote}
This year prices have gone up a lot.
\end{quote}

Built on \ta{விலை}, a price.

\begin{cousinweb}
\textbf{\ta{விலைவாசி}} (\emph{vilaivāci}) is built on \ta{விலை}, \textquotedblleft{}a price\textquotedblright{} (DEDR 5421); the origin of its second part is not settled.
\end{cousinweb}

\begin{grammarlens}[title={What you've built}]
One more word for getting around.
\end{grammarlens}

\subsection*{Your turn}
\begin{itemize}
\raggedright
  \item \emph{Say it:} \emph{vilaivāci}
  \item \emph{Say it:} \emph{vilaivāci}, once more
  \item \emph{Say it:} \emph{inta āṇṭu vilaivāci mikavum ēṟiviṭṭatu.} — This year prices have gone up a lot.
  \item \emph{Recall:} say the Tamil for a fake, then the Tamil for luxury, then say \emph{vilaivāci} again
\end{itemize}

\begin{tcolorbox}[breakable,colback=teal!4,colframe=teal!35!black,title={Before you move on}]
What does \ta{விலைவாசி} mean? (\textquotedblleft{}The cost of living; prices\textquotedblright{}.) Say it once more.
\end{tcolorbox}
\section[\texorpdfstring{\ta{நோட்டு} (nōṭṭu)}{நோட்டு (nōṭṭu)}]{\ta{நோட்டு} (nōṭṭu) — a banknote}

\label{lesson:TA-C450-nottu}



\begin{quote}
Before the new one: say the Tamil for luxury, then the Tamil for the cost of living.
\end{quote}

\subsection*{You'll want to know: \ta{நோட்டு}}
\textbf{\ta{நோட்டு}} — \emph{nōṭṭu} — \textquotedblleft{}a banknote\textquotedblright{}.

\begin{quote}
\ta{அவர்} \ta{ஐநூறு} \ta{ரூபாய்} \ta{நோட்டைக்} \ta{கொடுத்தார்}.
\end{quote}

\begin{quote}
\emph{avar ainūṟu rūpāy nōṭṭaik koṭuttār.}
\end{quote}

\begin{quote}
He gave a five-hundred-rupee note.
\end{quote}

Before -\ta{ஐ} the final -\ta{உ} drops: \ta{நோட்டை}, the note (object).

\begin{cousinweb}
\textbf{\ta{நோட்டு}} (\emph{nōṭṭu}) is a loan from English \emph{note}, with the doubled stop and final -\ta{உ} Tamil gives loans.
\end{cousinweb}

\begin{grammarlens}[title={What you've built}]
One more word for getting around.
\end{grammarlens}

\subsection*{Your turn}
\begin{itemize}
\raggedright
  \item \emph{Say it:} \emph{nōṭṭu}
  \item \emph{Say it:} \emph{nōṭṭu}, once more
  \item \emph{Say it:} \emph{avar ainūṟu rūpāy nōṭṭaik koṭuttār.} — He gave a five-hundred-rupee note.
  \item \emph{Recall:} say the Tamil for luxury, then the Tamil for the cost of living, then say \emph{nōṭṭu} again
\end{itemize}

\begin{tcolorbox}[breakable,colback=teal!4,colframe=teal!35!black,title={Before you move on}]
What does \ta{நோட்டு} mean? (\textquotedblleft{}A banknote\textquotedblright{}.) Say it once more.
\end{tcolorbox}
